% Imporditud: tools/import_eksperimendid.py
% Allikas: Eesti füüsikaolümpiaadi eksperimendi ülesannete
%   kogu (Rael Kalda, 2023), ülesanne Ü111, lahendus L111.
\setAuthor{Erkki Tempel}
\setRound{piirkonnavoor}
\setYear{2015}
\setNumber{G E2}
\setDifficulty{3}
\setTopic{geomeetriline optika}
\setVahendid{Nõguslääts, mõõtejoonlaud, ruuduline paber}
\setPunktid{12}
\setAllikas{Eksperimendi kogu Ü111, lahendus lk 87}
\setLahendusseis{toores}

\prob{Läätse fookuskaugus}
Leidke nõgusläätse fookuskaugus 20 \% täpsusega.

\hint

\solu
Vaatame ruudulist paberit läbi läätse ning võrdleme ruutude suurusi, mida näeme otse paberit vaadates ning läbi läätse vaadates (vt. joonis).

2 ruutu

4 ruutu

F

Kuna me peame leidma läätse fookuskauguse 20 \% täpsusega, siis võime teha lihtsustuse, et kaugelt vaadates on kujutis objekti asukohas. Kõige lihtsam on läätse fookuskaugust leida siis, kui nähtav kujutis on kaks korda väiksem kui ese. Sellisel juhul ühtib läätse fookuskaugus eseme kaugusega läätsest (vt joonis). Läätse fookuskaugus on sellise meetodiga mõõtes 16 cm $\pm$ 1 cm.

Hindamine: Idee võrrelda eseme suurust kujutise suurusega --- 3 p. Lihtsustus, et kujutis on objekti asukohas --- 3 p. Fookuskauguse leidmine eseme ja kujutise suuruste suhte kaudu --- 4 p. Kordusmõõtmised (3 mõõtmist) --- 2 p.
\probend
